\chapterrule
\section{位置编码、缩放与多头注意力的手算}
\label{intro:22a-position-multihead}
\chaptercompass{交换 Token 时，模型怎么知道位置？}{对照位置向量、缩放点积和两个注意力头的独立读取。}{$X:T\times d\to QK^\top:T\times T\to Y:T\times d$}{位置表可学习；Q/K/V 与输出投影由训练确定}
\modelpipelinefocus{3}{位置与多头注意力}

% Keep this matrix-heavy calculation continuous without a separate exercise tail.
\begingroup
\setlength{\parskip}{2pt}
\setlength{\abovedisplayskip}{4pt}
\setlength{\belowdisplayskip}{4pt}
\setlength{\abovedisplayshortskip}{2pt}
\setlength{\belowdisplayshortskip}{2pt}
\subsection{交换两个 Token：统一三维参数为什么需要位置？}

前面已经确定唯一三维 Embedding。为了比较顺序，
这里只拿两个已出现的 Token“春天、花”：
\[
X=\begin{bmatrix}1.2&0&-0.2\\-0.3&1.1&0.2\end{bmatrix}
\in\R^{2\times3},\qquad
X'=\begin{bmatrix}-0.3&1.1&0.2\\1.2&0&-0.2\end{bmatrix}.
\]
保持与统一教学参考脚本相同的 $W_Q,W_K,W_V$，先做\textbf{无位置、
无因果 Mask} 的一次双向 Attention。这里只隔离位置对比，
暂不执行 LayerNorm 与残差：
\[
O(X)=\softmax\left(\frac{(XW_Q)(XW_K)^\top}{\sqrt3}\right)XW_V
\approx\begin{bmatrix}
0.6693&0.4934&0.0066\\0.2019&0.6886&0.1683
\end{bmatrix}.
\]
交换输入两行后，$O(X')$ 的第一行就是上式第二行，第二行就是上式第一行。
这叫行置换的\textbf{等变性}：输出随输入换行，却没有额外标记“原来在哪个位置”。

\subsection{加入位置绑定的向量，再观察交换后的区别}
这次位置消融对照使用\textbf{同一组给定的位置数值}；
在实际训练中，本书三维主例的位置向量表 $P_{\rm pos}$ 仍是\textbf{可训练参数}。
也可采用正余弦位置编码。

位置向量取自贯通主实验的\textbf{前两行}，不因交换 Token 而移动：
\[
{P_{\rm pos}}_{1:2}=\begin{bmatrix}
0.10&0.05&-0.10\\0&-0.05&0.12
\end{bmatrix},\quad
X_+=X+{P_{\rm pos}}_{1:2},\quad X'_+=X'+{P_{\rm pos}}_{1:2}.
\]
仍用\textbf{相同的三组可训练投影}，只是输入加了位置信息：
\begin{equation}
O(X_+)\approx\begin{bmatrix}
0.7680&0.5138&-0.0265\\0.2324&0.6726&0.2092
\end{bmatrix},\qquad
O(X'_+)\approx\begin{bmatrix}
0.2812&0.7145&0.1440\\0.6956&0.4854&0.0441
\end{bmatrix}.
\label{eq:pos-swap-d3}
\end{equation}
注意交换后的第一行\textbf{并不}等于交换前的第二行：
例如第一坐标分别为 $0.2812$ 和 $0.2324$。
位置向量提供了独立于内容的位置信号。
两种顺序下的输出也没有退化成两行完全相同。

\textbf{位置消融条件：}这里暂不加因果 Mask、LayerNorm 和残差；
完整自回归 Block 则包含这些步骤。

\subsection{\texorpdfstring{为什么点积要除以 $\sqrt{d_k}$？算一次波动尺度}{为什么点积要除以 Key 维度平方根？}}
\phantomsection\label{concept:scaled-dot-product}

\noindent\emph{数学随手查：}独立零均值随机量相加时，方差相加，标准差取方差的平方根。
求和与平方根可回看\hyperref[app:math-tools]{数学速查}。


假设两个 Key/Query 向量的每一维都是\textbf{独立、等概率取 $\pm1$}
的随机变量，且维度之间也独立。
每个分量乘积平均为 $0$、方差为 $1$。
$d_k$ 个分量求和，点积的方差为 $d_k$，
标准差为 $\sqrt{d_k}$；当 $d_k=64$ 时约为 $8$。

比如两个候选的未缩放分数恰好是 $(8,0)$，
Softmax 第一项为
$\frac{e^8}{e^8+1}\approx0.9997$，
另一项几乎为零。除以 $\sqrt{64}=8$ 后，
分数变成 $(1,0)$，第一项变成
$\frac e{e+1}\approx0.731$。
缩放能防止高维点积的典型量级随着维度扩大而直接变得很尖锐；
这不是保证所有 Attention 都不会饱和的定理。

\subsection{两个头各算二维，再拼接：数字完整走一遍}
\phantomsection\label{concept:multihead}

\localfixture{四维双头}
取 $d_{\rm model}=4$，每头宽度为 $2$，观察最后位置读取两个 Key。
给定两头的最后位置 Query 为
$q^{(1)}=(0,1)$、$q^{(2)}=(1,0)$，两头的 Key 都为 $I_2$，
Value 分别为
\[
V^{(1)}=\begin{bmatrix}2&0\\0&2\end{bmatrix},
\qquad
V^{(2)}=\begin{bmatrix}0&1\\1&0\end{bmatrix}.
\]
按“点积 $\to$ 除以 $\sqrt2$ $\to$ Softmax $\to$ 加权 Value”逐行计算：
\[
\begin{aligned}
s^{(1)}&=(0,1)/\sqrt2,&
a^{(1)}&\approx(0.330,0.670),\\
o^{(1)}&=a^{(1)}V^{(1)}=(0.660,1.340),\\
s^{(2)}&=(1,0)/\sqrt2,&
a^{(2)}&\approx(0.670,0.330),\\
o^{(2)}&=a^{(2)}V^{(2)}=(0.330,0.670).
\end{aligned}
\]
若省略缩放，第一头权重会变成 $(0.269,0.731)$；
本例统一使用上面的缩放结果。
沿\textbf{特征列}拼接，得到
$O_{\rm cat}=(0.660,1.340,0.330,0.670):1\times4$。
再取演示投影
\[
W_O=\begin{bmatrix}1&0&1&0\\0&1&0&1\end{bmatrix}^{\top},
\qquad O_{\rm cat}W_O=(0.990,2.010):1\times2.
\]
\textbf{多头拼接各头读出的特征，再投影到目标宽度。}
这里用 $4\times2$ 演示投影；完整四维残差 Block 的输出投影则应为 $4\times4$。

\chapterchecktasks{22a-position-multihead}{正文基础三题}{
\item \textbf{预测：}句子里同一个 Token 在不同位置出现，只有内容 Embedding 而没有任何位置机制时，网络失去什么信息？为什么需要位置表示？\chapterinlinehint{交换两个相同 Token 的位置，检查输入有无变化。}
\item \textbf{形状：}两个头各输出 $3\times2$，按特征轴拼接后得什么形状？再用 $W_O:4\times6$ 投影得到什么形状？\chapterinlinehint{先累加拼接宽度，再与输出投影核对内侧维度。}
\item \textbf{找错：}同学认为“Attention 分数除以 $\sqrt{d_k}$ 后，所有相似度都必须小于 $1$”。错在哪里？\chapterinlinehint{比较尺度变换与区间裁剪的定义。}
}
\chapterrecallanswer{位置表示携带序列顺序，多头允许独立匹配视角；缩放用于控制点积幅度，不是把分数截成概率。}



\begin{conclusion}
\textbf{位置表示}标记顺序，\textbf{多头}在不同子空间读取。
下一章用 Norm、残差与 FFN 把它们接进完整 Block。
\end{conclusion}
\endgroup
